% ============================================================
% fig_hopfield_energia.tex — Paisaje de energía de Hopfield
% ============================================================
\begin{figure}[H]
\centering
\begin{tikzpicture}[scale=1.1]

% Ejes
\draw[-{Stealth}, thick] (-0.3, 0) -- (7.2, 0) node[right, font=\small] {espacio de estados};
\draw[-{Stealth}, thick] (0, -0.3) -- (0, 4.0) node[above, font=\small] {$E$};

% Curva de energía con varios mínimos locales
% Fix: usar grados directamente (sin sufijo r) con factores escalados
\draw[blue!70, very thick, domain=0.2:7.0, samples=200]
    plot (\x, {
        0.6*sin(51.43*\x) +
        0.4*sin(102.86*\x) +
        0.3*cos(77.14*\x) +
        1.8
    });

% Mínimos (patrones almacenados)
\fill[green!60!black] (1.05, 0.82) circle (3.5pt);
\fill[green!60!black] (3.55, 0.62) circle (3.5pt);
\fill[green!60!black] (5.90, 0.95) circle (3.5pt);

\node[green!50!black, font=\footnotesize, above right] at (1.05, 0.82) {patrón $\xi^1$};
\node[green!50!black, font=\footnotesize, above right] at (3.55, 0.62) {patrón $\xi^2$};
\node[green!50!black, font=\footnotesize, above right] at (5.90, 0.95) {patrón $\xi^3$};

% Estado inicial ruidoso → convergencia hacia mínimo
\fill[red!70]    (2.3, 2.85) circle (3pt) node[above, font=\footnotesize, red!70] {entrada ruidosa};
\draw[-{Stealth}, red!60, thick, dashed]
    (2.3, 2.85) .. controls (2.1, 2.0) and (1.7, 1.4) .. (1.05, 0.85);
\node[red!60, font=\tiny] at (1.5, 1.9) {convergencia};

% Etiqueta eje Y
\node[font=\footnotesize, gray, rotate=90] at (-0.55, 2.0) {Energía $E$};

\end{tikzpicture}
\caption{Paisaje de energía de una red de Hopfield. Cada mínimo local corresponde
a un patrón almacenado $\xi^k$. Ante una entrada ruidosa o incompleta, la dinámica
de la red desciende por el gradiente de energía hasta converger al mínimo más
cercano, recuperando el patrón original.}
\label{fig:hopfield_energia}
\end{figure}
